\begin{claim}{F.32}{Proved}{Finite affine marginal image equality when the first character is one and the boundary character is nontrivial}
Use the original bases $Y_i^-=X_i$ and
$Y_i^+=\mathcal B_+(X_i)$ of F.23.13.
For every finite field $k$ and continuous
$\rho:V_\pm\to\operatorname{Aff}_1(k)$ whose
scalar character belongs to the exact locus
\[
 \chi(Y_1^\pm)=1,\qquad
                 h=\prod_{i=1}^7\chi(Y_i^\pm)\ne1,
                                                               \tag{F.32.1}
\]
the actual marginal images satisfy
\[
                         \rho(\mathcal K_\pm)=\rho(R_\pm).
                                                               \tag{F.32.2}
\]
All finite characteristics, scalar-field
extension degrees and larger original
coefficient fields are included.

Put $C=\chi(R_\pm)$,
$k_1=\mathbb F_p(C)$ and
$k_0=\mathbb F_p(\chi(V_\pm))$.
For $C\ne\{1\}$, let $q_1=|k_1|$,
$m=[k_0:k_1]$ and
$\chi_j=\chi^{q_1^j}$, $0\leq j<m$.
The summed continuous $H^1$ restriction
for these characters is injective.
The corresponding cocycle restriction
has exactly the principal tuples with
parameter sum zero as its kernel,
and has rank $6m+1$ on $7m$ coordinates.
\end{claim}

\begin{proof}
\emph{The actual scalar fields and coefficient maps.}
Work first with the original negative
marginal $K=\mathcal K_-\leq R=R_-\leq V=V_-$.
Write its character in the unique form
\[
 \begin{gathered}
 \chi=(1,s,t,u,v,w,z),\qquad h=stuvwz\ne1,\\
 C=\chi(K)=\chi(R)=\langle u,v,w,z\rangle,\qquad
 k_1=\mathbb F_p(u,v,w,z),\\
 k_0=k_1(s,t)=k_1(h,t),\qquad
                       s=\frac{h}{tuv wz}.
 \end{gathered}                                                \tag{F.32.3}
\]
Every displayed scalar is nonzero.
The equality of scalar images is the
image of the actual normal-generation
equality
$\overline{\langle\!\langle K\rangle\!\rangle}_R=R$
from F.21 and the marginal map of F.22
under the abelian scalar homomorphism.
The deleted generators $X_4,\ldots,X_7$
give exactly the four generators of $C$.
The identities for $h$ and $s$ prove
both inclusions in the displayed
equality of fields.

For $C=\{1\}$, the image of $R$ is
an abelian translation group, so this
same normal-generation map proves
(F.32.2). For $C\ne\{1\}$ and $m=1$,
F.25 gives (F.32.2), and F.24 gives
the asserted individual restriction
injection and cocycle kernel.
We henceforth take $C\ne\{1\}$ and $m>1$.

The original Laurent coefficient map is
\[
 \begin{aligned}
 \mathbb Z[t_1^{\pm1},\ldots,t_7^{\pm1}]
   &\longrightarrow
       A=\mathbb Z[h^{\pm1},t^{\pm1},u^{\pm1},v^{\pm1},w^{\pm1},z^{\pm1}],\\
 (t_1,t_2,t_3,t_4,t_5,t_6,t_7)
   &\longmapsto (1,h/(tuv wz),t,u,v,w,z).
 \end{aligned}                                                \tag{F.32.4}
\]
After the first character is set to one,
the change from the independent variables
$(s,t,u,v,w,z)$ to $(h,t,u,v,w,z)$ is
the exact Laurent-ring isomorphism
\[
 \begin{gathered}
 s^at^bu^cv^dw^ez^f
    \longmapsto h^at^{b-a}u^{c-a}v^{d-a}w^{e-a}z^{f-a},\\
 h^at^bu^cv^dw^ez^f
    \longmapsto s^at^{b+a}u^{c+a}v^{d+a}w^{e+a}z^{f+a}.
 \end{gathered}                                                \tag{F.32.5}
\]
Both composites fix every monomial.
Adding exponent vectors proves that
both maps preserve multiplication;
extension by integer linearity proves
the asserted inverse ring maps.
In $A$ the letter $s$ will always denote
the specific Laurent monomial in (F.32.3).

\emph{Boundary coordinates for the actual principal quotient.}
The continuous cocycle space for $\chi$
is $k_0^7$, in original generator order.
Indeed every such tuple defines a
homomorphism of the discrete free group
to the finite affine group, and that
homomorphism has a unique continuous
extension to $V=\widehat F_7$.
Its principal line is $k_0\delta$, where
\[
 \begin{gathered}
 \delta=(0,s-1,t-1,u-1,v-1,w-1,z-1)^{\mathsf T},\\
 D=(1,1,s,st,stu,stuv,stuvw),\qquad
                            D\delta=h-1,\\
 N=I_7-\frac{\delta D}{h-1},\qquad H=\ker D.
 \end{gathered}                                                \tag{F.32.6}
\]
The row $D$ evaluates a cocycle on
the original word $P_0=X_1\cdots X_7$.
The nonzero scalar $h-1$ makes
$\delta$ nonzero and defines $N$.
Direct multiplication gives
$N\delta=0$, $DN=0$ and $N^2=N$.
For every $a\in k_0^7$ one has
\[
       a=Na+\delta\frac{Da}{h-1},\qquad
       \ker N=k_0\delta,\qquad \operatorname{im}N=H.
                                                               \tag{F.32.7}
\]
Consequently $[a]\mapsto Na$ and
$b\mapsto[b]$, $b\in H$, are mutually
inverse maps from the actual
$Q_\chi=H^1_{\rm cont}(V,k_{0,\chi})$
to $H$. This proves the precise
principal-quotient identification.

Let $E$ take the last six coordinates,
and let $J:k_0^6\to k_0^7$ have columns
\[
 \begin{gathered}
 J=(e_2-e_1,e_3-se_1,e_4-ste_1,
                         e_5-stue_1,e_6-stuve_1,e_7-stuvwe_1),\\
 DJ=0,\qquad EJ=I_6,\qquad JE=I_7-e_1D,\\
 ENJ=I_6,\qquad JEN=N.
 \end{gathered}                                                \tag{F.32.8}
\]
Thus $J$ and $E|_H$ are explicit
inverse coordinate maps on $H$.
The actual principal quotient in six
coordinates is $EN$, with section $J$.
These coordinates have the original
order $X_2,\ldots,X_7$.

Retain each of the eleven actual
F.23 source automorphisms
$B_l=\operatorname{Ad}(u_l)\alpha_l$.
Here and below $0\leq l\leq10$ is
the original zero-based source index:
$0$ is \texttt{old0}, $1,\ldots,5$
are \texttt{old\_1\_0},\ldots,\texttt{old\_1\_4},
and $6,\ldots,10$ are
\texttt{old\_-1\_0},\ldots,\texttt{old\_-1\_4}.
Its position in a one-based list is
always $l+1$.
Both $B_l$ and its inverse preserve
the actual subgroup $K$.
The pure Artin part preserves every
abelian character and fixes $P_0$.

Let $A_l$ be the seven-coordinate
cocycle action of $B_l$, $M_l$ that
of $\alpha_l$, $a_l$ the row evaluating
$d(u_l)$, and $\lambda_l=\chi(u_l)$.
The actual cocycle formulas give
\[
 \begin{gathered}
 A_l=\lambda_lM_l-\delta a_l,\qquad
 M_l\delta=\delta,\qquad DM_l=D,\\
 a_l\delta=\lambda_l-1,\qquad A_l\delta=\delta,\\
 U_l=A_l+\delta a_l=\lambda_lM_l,
             \qquad NA_l=U_lN,\\
 T_l=EU_lJ,\qquad U_lJ=JT_l.
 \end{gathered}                                                \tag{F.32.9}
\]
For the middle intertwiner, $N\delta=0$
cancels the displayed decoration term,
and $M_l\delta=\delta$, $DM_l=D$
give $NM_l=M_lN$.
Thus $T_l$ is the exact source action
on the actual quotient (F.32.8).
Its entries belong to $A$; the only
division in identifying that quotient
was the already nonzero $h-1$.

For an entirely integral check put
$B=(h-1)I_7-\delta D$. Then
\[
 \begin{gathered}
 B\delta=0,\quad DB=0,\quad B^2=(h-1)B,\quad
 BA_l=U_lB,\\
 DU_l=\lambda_lD,\qquad U_l\delta=\lambda_l\delta,\\
 (\lambda_0,\ldots,\lambda_{10})=
 (1,1,t,1,1,1,u^{-2},(ht)^{-1},v^{-2},w^{-2},z^{-2}).
 \end{gathered}                                                \tag{F.32.10}
\]
Every scalar in the last line is
evaluated from the decoration of the
same labelled original word.

\emph{The invariant five-dimensional kernel and its quotient.}
Because $\chi(X_1)=1$, evaluation
$d\mapsto d(X_1)$ vanishes on the
principal line. On the coordinates
(F.32.8) its negative is the row
$b=(1,s,st,stu,stuv,stuvw)$.
Set $W_\chi=\ker b$ and use the exact
maps and basis
\[
 \begin{gathered}
 \iota(x_1,\ldots,x_5)=
 (-sx_1-stx_2-stux_3-stuvx_4-stuvwx_5,x_1,\ldots,x_5),\\
 \pi(\zeta)=b\zeta=-d(X_1),\qquad \sigma(a)=ae_1,\qquad
 \pi\sigma=I_1,\\
 \zeta=\iota(\zeta_2,\ldots,\zeta_6)+\sigma\pi(\zeta),\\
 F=(e_2-se_1,e_3-ste_1,e_4-stue_1,
                         e_5-stuve_1,e_6-stuvwe_1,e_1),\\
 F^{-1}(\zeta)=(\zeta_2,\zeta_3,\zeta_4,\zeta_5,\zeta_6,b\zeta).
 \end{gathered}                                                \tag{F.32.11}
\]
Both composites of the displayed basis
maps are identity. In particular no
additional scalar is inverted here.
Direct evaluation of the original words
gives $d(B_l(X_1))=\nu_l d(X_1)$ and
\[
 \begin{gathered}
 S_l=F^{-1}T_lF=
       \begin{pmatrix}G_l&*\\0&\nu_l\end{pmatrix},\qquad
                              G_0=h^{-1}I_5,\\
 (\nu_0,\ldots,\nu_{10})=
 (1,u^{-2},t^{-1},v^{-2},w^{-2},z^{-2},1,t/h,1,1,1).
 \end{gathered}                                                \tag{F.32.12}
\]
This proves invariance of $W_\chi$
for every actual source, and identifies
its one-dimensional quotient by the
actual map $\pi$.
The integral certificate checks the
full original first row of $A_l$,
the full block identity, and the cleared
seven-coordinate identity
$BA_lJF=(h-1)JFS_l$.

\emph{Eight moving-root identities and explicit evaluation rows.}
For each row of the following table
let $Q_l(X)$ be the product of $X-a$
over all listed base roots $a$, with
their displayed multiplicities.
Let $b_l$ be its moving root.
\[
\begin{array}{c|c|c}
 l&\text{base roots of }Q_l&b_l\\\hline
 1&1,(vw)^{-1}&(tz)^{-1}\\
 3&1,(uz)^{-1}&(tw)^{-1}\\
 4&1,(vz)^{-1}&(tu)^{-1}\\
 5&1,(uw)^{-1}&(tv)^{-1}\\
 6&1,(vw)^{-1},(h/u)^2&(tz)^{-1}\\
 8&1,(uz)^{-1},(h/v)^2&(tw)^{-1}\\
 9&1,(vz)^{-1},(h/w)^2&(tu)^{-1}\\
 10&1,(uw)^{-1},(h/z)^2&(tv)^{-1}
\end{array}.
                                                               \tag{F.32.13}
\]
The following identities hold over $A$,
before any specialization or cancellation:
\[
 \begin{gathered}
 Q_l(G_l)(G_l-b_lI_5)=0,\qquad
 q_lG_l=b_lq_l,\\
 q_lQ_l(G_l)=Q_l(b_l)q_l,\qquad
 q_l=\widehat q_l\Delta_t,\qquad
               \Delta_t=\operatorname{diag}(1,t-1,t-1,t-1,t-1).
 \end{gathered}                                                \tag{F.32.14}
\]
The diagonal identity defines $q_l$
without dividing by $t-1$.
Here are all eight rows explicitly.
The first four are
\[
 \begin{gathered}
 \widehat q_1^{\mathsf T}=\frac1{uv wz}
 \begin{pmatrix}
 -h(z-1)\\-h(z-1)\\-hu(z-1)\\-huv(z-1)\\-uv w(hz-h-z)
 \end{pmatrix},\qquad
 \widehat q_8^{\mathsf T}=\frac1{uv w}
 \begin{pmatrix}
 -h(w-1)\\-h(w-1)\\-hu(w-1)\\-uv(hw-h-w)\\-uv w(h-1)(w-1)
 \end{pmatrix},\\
 \widehat q_9^{\mathsf T}=\frac1{uv wz}
 \begin{pmatrix}
 h(u-1)\\h(u-1)-uv wz\\hu(u-1)\\huv(u-1)\\huv w(u-1)
 \end{pmatrix},\qquad
 \widehat q_{10}^{\mathsf T}=
                 \begin{pmatrix}1-v\\1-v\\u\\0\\0\end{pmatrix}.
 \end{gathered}                                                \tag{F.32.15}
\]
The other four use the explicit
polynomials
\[
 \begin{aligned}
 P_4={}&h^2\{u(vw+wz-w-z)-vw\}
                  +hu(-vw-2wz+w+z)+uwz,\\
 P_6={}&(z-1)\{h^2[vw(u-1)+uw(z-1)-uz]
                         +huz(1-2w)\}+uwz^2.
 \end{aligned}                                                \tag{F.32.16}
\]
Their rows are
\[
 \begin{gathered}
 \widehat q_3^{\mathsf T}=\frac1{uwz}
 \begin{pmatrix}
 -h(w-1)\\-(h+uwz)(w-1)\\0\\hu(uw-w+1)\\huw(u-1)(w-1)
 \end{pmatrix},\\
 \widehat q_4^{\mathsf T}=\frac1{u^2wz}
 \begin{pmatrix}
 h^2(u-1)\\
 h^2(u-1)+huv wz(u-1)-u^2v wz\\
 -hu(u-1)[h(u-1)-u]\\
 -hu(u-1)\{h[(u-1)v+u(z-1)]-u(v+z-1)\}\\
 -u(u-1)P_4
 \end{pmatrix},\\
 \widehat q_5^{\mathsf T}=\frac1{uwz}
 \begin{pmatrix}
 h(v-1)\\h(v-1)\\-hu[hz+(v-1)(z-1)]\\
 -u(v-1)[h^2z+hv(z-1)+h(w-1)z-wz]\\
 -uw(v-1)[h^2z+hv(z-1)+h(w-2)z-(v+w-1)z]
 \end{pmatrix},\\
 \widehat q_6^{\mathsf T}=\frac1{u^2wz^2}
 \begin{pmatrix}
 h^2(z-1)\\h(h+uv wz)(z-1)\\-h^2u(u-1)(z-1)\\
 -hu(z-1)\{h[u(v+z-1)-v]-uz\}\\-uP_6
 \end{pmatrix}.
 \end{gathered}                                                \tag{F.32.17}
\]
Every denominator in these displays
is a Laurent monomial. The direct
left-eigenrow identities in (F.32.14)
are checked from the actual matrices;
they do not require a rank assertion
to persist under specialization.
The last polynomial evaluation in
(F.32.14) also follows by multiplying
the direct left-eigenrow identity
through each factor of $Q_l$.

For $a=3,4,5,6$, let $\widehat R_a$
be the square matrix with rows
$\widehat q_1,\widehat q_8,\widehat q_9,
\widehat q_{10},\widehat q_a$, in that order.
Let $R_a$ use the corresponding actual
rows $q_l$. The exact determinants are
\[
 \begin{gathered}
 R_a=\widehat R_a\Delta_t,\qquad
                 \det R_a=(t-1)^4\det\widehat R_a,\\
 \begin{array}{c|c}
 a&\det\widehat R_a\\\hline
 3&h u(h-1)(w-1)/(wz)\\
 4&-h(h-1)(u-1)(hw-h-w)/w^2\\
 5&-h(h-1)(v-1)\\
 6&-h(h-1)(z-1)(hw-h-w)/(w^2z)
 \end{array}.
 \end{gathered}                                                \tag{F.32.18}
\]
This is a column-scaling identity
over $A$, including $t=1$.
Its use as an invertibility statement
will occur only after the relevant
field argument proves $t\ne1$.

\emph{Comparison of core subspaces in distinct Frobenius factors.}
Let $\sigma:k_0\to k_0$ be the actual
$k_1$-Frobenius $a\mapsto a^{q_1}$.
Its order is $m$. Put
$h_j=\sigma^j(h)$ and $t_j=\sigma^j(t)$;
the four deleted scalars are fixed.
The coordinates and matrices above
specialize to $Q_j=Q_{\chi_j}$ and
$W_j\subset Q_j$.
For each $j$ the two-way field equality
in (F.32.3), after this automorphism,
gives
\[
                       k_0=k_1(h_j,t_j).
                                                               \tag{F.32.19}
\]

Suppose that nonzero invariant subspaces
$S_j\subset W_j$ and $S_{j'}\subset W_{j'}$,
with $j\ne j'$, have a linear isomorphism
intertwining all eleven sources.
Source zero is $h_j^{-1}I_5$ and
$h_{j'}^{-1}I_5$, respectively, so
$h_j=h_{j'}$. Use the actual automorphism
$\tau=\sigma^{j'-j}$, with negative powers
interpreted by the inverse automorphism,
and put
\[
 L=\{a\in k_0:\tau(a)=a\},\qquad
       k_1\subset L,\quad h_j\in L,\quad
                         t_j\notin L.
                                                               \tag{F.32.20}
\]
Here $L$ is a field because $\tau$
preserves the field operations and
inverses. It is proper, since
$0<|j'-j|<m$ and $\sigma$ has order $m$.
If $t_j$ belonged to $L$, (F.32.19)
would put all of $k_0$ in $L$, a
contradiction. The same reasoning shows
$t_{j'}\ne t_j$: equality would make
$\tau$ fix both generators in (F.32.19).
Also $t_{j'}=\tau(t_j)$ cannot belong
to $L$, since applying $\tau^{-1}$ to
an element of $L$ fixes that element.

For every source in (F.32.13), all base
roots lie in $L$ and agree for the
two factors. Its moving roots are
$\eta/t_j$ and $\eta/t_{j'}$ for one
nonzero $\eta\in k_1$. They are distinct,
and neither belongs to $L$; otherwise
their inverses multiplied by $\eta$
would put $t_j$ or $t_{j'}$ in $L$.
Consequently each factor of $Q_l(b_{l,j})$
is nonzero, including any repeated
base-root factor.

Transport the source action on $S_{j'}$
back to $S_j$ by the given isomorphism.
For a fixed source it is the same
operator on $S_j$, and both annihilators
in (F.32.14) apply to it.
Subtracting these two polynomial
identities gives
$(b_{l,j'}-b_{l,j})Q_l(G_{l,j})=0$
on $S_j$, so $Q_l(G_{l,j})$ vanishes
there. The direct row evaluation in
(F.32.14) and the proved nonzero
$Q_l(b_{l,j})$ now give
$q_{l,j}|_{S_j}=0$ for all eight rows.
This deduction uses only exact
intertwiners and polynomial identities.

At this comparison $t_j\notin L$,
so $t_j-1\ne0$; also $h_j-1\ne0$
by (F.32.1). Every displayed monomial
factor is nonzero. The determinant
cover (F.32.18) is exhaustive as follows.
For $w\ne1$ use $a=3$.
For $w=1$ and $v\ne1$ use $a=5$.
For $w=v=1$ and $u\ne1$ use $a=4$:
the remaining factor $h_jw-h_j-w$
is exactly $-1$.
For $w=v=u=1$, nontriviality of $C$
forces $z\ne1$, and $a=6$ has the
same factor $-1$.
Each case provides a nonzero actual
determinant $\det R_a$.
Thus the eight rows cannot all vanish
on a nonzero subspace, a contradiction.
There are no isomorphic nonzero
invariant core subspaces across
distinct Frobenius factors.
Every argument and every nonzero
factor holds also in characteristic two.

\emph{All other simple-subspace types.}
A simple invariant subspace of $Q_j$
either lies in $W_j$, or intersects
$W_j$ trivially. In the latter case
the actual quotient map $\pi$ restricts
to an injection into its one-dimensional
quotient. Its image is the entire
one-dimensional space, so it has
exactly the character $\nu_j$ of
(F.32.12), with the inverse of this
restricted quotient as the inverse map.
This proves the classification of
possible simple-subspace types used
here without a semisimplicity assumption.

A core subspace from any factor cannot
be isomorphic to such a quotient
character from any factor: source zero
is respectively $h_j^{-1}$ and $1$,
which are unequal because $h_j\ne1$.
For two quotient characters, equality
at source $2$ gives $t_j=t_l$, and
equality at source $7$ then gives
$h_j=h_l$. Equation (F.32.19) makes
$\sigma^{l-j}$ the identity on $k_0$,
which forces $j=l$ in the indicated
index range. Inverse intertwiners cover
both orientations of each comparison.
Together with the core comparison,
these are all three unordered type
pairs: core/core, core/quotient and
quotient/quotient.

\emph{Summed restriction and the exact affine consequence.}
Every $\chi_j$ restricts to the same
character on $K$, because its image
there is $C\subset k_1^\times$.
There is consequently the actual map
\[
 \Phi:\bigoplus_{j=0}^{m-1}Q_j
       \longrightarrow H^1_{\rm cont}(K,k_{0,\chi|K}),\qquad
                   ([d_j])_j\longmapsto\sum_j[d_j|K].
                                                               \tag{F.32.21}
\]
The same eleven actual source
automorphisms act on both sides, and
restriction commutes with them because
they preserve $K$ and the characters.
Each individual restriction is injective
by F.24 and $C\ne\{1\}$.

A nonzero kernel of $\Phi$ would
contain a nonzero invariant subspace
$T$ of least dimension. Each nonzero
coordinate projection $p_j|_T$ is
injective: its kernel is invariant
and is a proper subspace of $T$.
Its image is simple, because the
inverse projection pulls any proper
nonzero invariant subspace back to one
in $T$. Two nonzero projections give
the explicit inverse intertwiners
\[
 \begin{gathered}
 p_l\circ(p_j|_T)^{-1}:p_j(T)\longrightarrow p_l(T),\\
 p_j\circ(p_l|_T)^{-1}:p_l(T)\longrightarrow p_j(T).
 \end{gathered}                                                \tag{F.32.22}
\]
The exhaustive three comparisons exclude
these for $j\ne l$. Exactly one
projection is therefore nonzero, and
$T$ lies in a single factor. There
$\Phi$ is the injective individual
restriction, giving the final
contradiction. This proves the summed
$H^1$ injection.

Zero summed cocycle restriction now
makes every cocycle principal.
Their common restricted character has
nontrivial image $C$, so the exact
kernel on cocycles is
\[
 \left\{\bigl(b_j(\chi_j-1)\bigr)_{j=0}^{m-1}:
                              \sum_{j=0}^{m-1}b_j=0\right\}.
                                                               \tag{F.32.23}
\]
Conversely each tuple in this display
restricts with sum zero. Each principal
line has a nonzero generating vector;
the displayed kernel has dimension
$m-1$. The cocycle domain has its
original $7m$ coordinates, so its
restriction rank is $6m+1$.

Apply this statement to the actual
Frobenius coefficient tuple
$(d^{q_1^j})_{j=0}^{m-1}$. It gives
the exact trace-kernel identity
\[
 \operatorname{Tr}_{k_0/k_1}(d|K)=0
 \quad\Longleftrightarrow\quad
 d=b(\chi-1),\qquad
                      \operatorname{Tr}_{k_0/k_1}(b)=0.
                                                               \tag{F.32.24}
\]
For the forward direction, summed
restriction makes $d$ principal. Its
Frobenius powers have the unique
parameters $b^{q_1^j}$, whose sum is
the trace, and (F.32.23) gives the
stated vanishing. For the reverse
direction use $\chi(K)\subset k_1$
and the trace's $k_1$-linearity.

The explicit affine conjugation,
translation projection and trace-pairing
coefficient maps of F.26 take every
proper marginal affine image to a
nonprincipal element of this trace
kernel over the actual scalar field
$k_0$. Those maps preserve the scalar
character, both scalar fields, its
first coordinate and its boundary
product. They therefore retain the
exact locus (F.32.1).
Equation (F.32.24) excludes the
resulting nonprincipal cocycle, and
proves (F.32.2) for every original
finite coefficient field, including
fields larger than $k_0$.

\emph{Literal-word certificate, continuity and the positive marginal.}
The independent verifier imports none
of the four new producers and no
symbolic algebra package. It reconstructs
all eleven literal F.23 decorated
Artin words directly on the coefficient
chart (F.32.4), using the previously
reviewed integer Laurent word arithmetic.
At a positive crossing the meridian
and decoration pair is transformed by
$(x_j,x_{j+1};a_j,a_{j+1})\mapsto
(x_jx_{j+1}x_j^{-1},x_j;x_ja_{j+1},a_j)$;
at a negative crossing it is transformed
by $(x_{j+1},x_{j+1}^{-1}x_jx_{j+1};
a_{j+1},x_{j+1}^{-1}a_j)$.
The eighth meridian is evaluated at
the identity, and the final decoration
gives $u_l=a_8^{-1}$.
The verifier checks the label permutation,
the erased meridian, every preserved
character and the full boundary product.

Its $1088$ exact checks include the
full seven-coordinate actions and
decoration derivatives, the inverse
coefficient maps (F.32.5), all basis
inverses, the cleared intertwiners
(F.32.10)--(F.32.12), the eight direct
left-eigenrow identities and annihilators,
and all four actual determinant and
column-scaling identities.
They are integer Laurent identities,
so their stated field specializations
apply in every characteristic.
The nonmonomial division $h-1$ in
the principal quotient and all nonzero
factors used in the comparison have
been proved valid in the exact fields
before use. The evidence manifest
binds the original words, all coefficient
data, the programs and outputs, and
the completed dependencies.

All finite coefficient maps are
continuous. The original source
automorphisms are continuous by F.23;
the cocycles have the unique continuous
finite-affine extensions already described.
Their restrictions to the actual
closed subgroup $K$ and all displayed
coefficient compositions are therefore
continuous throughout the proof.
For the positive marginal use the
actual based isomorphism
$\mathcal B_+:V_-\to V_+$ of F.23.13
and its given inverse. It sends $X_i$
exactly to $Y_i^+$ and carries the
actual subgroups $\mathcal K_-,R_-$
onto their positive counterparts.
Precomposition by these inverse maps
preserves all seven character values,
the boundary product, both scalar
fields, and the maps in (F.32.21).
It therefore transports the proved
negative image equality and exact
kernel assertions to the positive
marginal. This proves the claim.
\end{proof}
